\section{Introduction}

Automatic short answer scoring (ASAS) is a critical task in educational assessment, where a model is trained to score student responses to open-ended questions.
Traditional approaches have largely relied on supervised learning from annotated answers to the same questions encountered at evaluation time, using either engineered features or neural representations \citep{burrowsErasTrendsAutomatic2015,riordanInvestigatingNeuralArchitectures2017}. As a result of supervised training, performance of these models often drops once confronted with unseen questions and domains.
Cross-question generalization, where a scorer must evaluate responses to questions absent from its training data, therefore remains an open challenge.
New questions introduce new target concepts, reference answers, and scoring criteria, which limit the transfer of patterns learned from seen questions \citep{dzikovskaSemEval2013Task2013,funayamaReducingCostCrossPrompt2023}.
Reducing this dependence on questions seen during training is critical for the practical deployment of ASAS systems, as it allows for the evaluation of student responses to new questions without requiring additional human annotation that is often not feasible in practical institutional contexts.

A large body of previous work tackles this challenge either by prompting large language models (LLMs) to score student answers~\citep{changAutomaticShortAnswer2024,chamiehLLMsShortAnswer2024} directly or by adapting the model architecture to better generalize to unseen questions~\citep{wangInjectRubricsShort2019,liSemanticFeatureWiseTransformation2021,funayamaReducingCostCrossPrompt2023,wuMatchingExemplarNext2023,sun-etal-2026-rubric}. There is also growing interest in using LLMs to synthesize training data for ASAS, to improve performance on unseen questions or address label imbalance \citep{yeZeroGenEfficientZeroshot2022,bexteLunchFreeOvercoming2025}. However, these works also show that student answers generated purely from LLM-prompting are less authentic and not always helpful for downstream training. 

Much existing work on data synthesis with LLMs for natural language understanding (NLU) has focused on general tasks such as sentiment analysis, intent detection, natural language inference (NLI), and topic classification~\citep{kumarDataAugmentationUsing2020,liSyntheticDataGeneration2023,zouFuseGenPLMFusion2024,ceginRoSERoundrobinSynthetic2025,pecherBetterGeneratorsClassifiers2026}, 
where labels usually carry a \textbf{fixed}, \textbf{input-independent} meaning that is often signaled by superficial lexical or semantic cues in the input text~\citep{joulinBagTricksEfficient2017,gururanganAnnotationArtifactsNatural2018,kaushikLearningDifferenceThat2020}.
ASAS differs from these tasks in two respects.
First, a useful synthetic answer must reflect a specific level of \textbf{understanding of the scientific concepts} involved, including partial knowledge and typical misconceptions, while also imitating the \textbf{linguistic style} of real students, such as informal phrasing and incomplete explanations.
Second, each score level's meaning is \textbf{question-specific}: it depends on the context and scoring criteria for each question rather than on fixed label semantics.
Established augmentation techniques that perturb seed examples from the training data via paraphrasing, insertion, or swapping~\citep{ceginLLMsVsEstablished2025} are therefore impractical, since surface-level edits cannot turn an answer to a seen question into an answer to a new one.

The goal of this paper is therefore to address these challenges by comprehensively investigating whether and how LLMs can synthesize realistic student responses to unseen open-ended science questions for ASAS training.

The contributions of this paper are as follows:
\begin{itemize}
    \item We propose a pipeline that uses existing annotated training data to synthesize new training data for out-of-distribution test data, e.g., unseen questions and domains. Since the pipeline requires only labeled examples, it can be readily applied to other NLU classification tasks.
    \item For this purpose, we investigate various ways of synthesizing realistic student responses to science questions for ASAS training, comparing few-shot prompting with synthesizers trained via supervised fine-tuning (SFT) and direct preference optimization (DPO) on annotated responses to seen questions. It turns out that a synthesizer trained on annotated responses to seen questions can readily generate realistic and useful synthetic responses to unseen questions, particularly when trained via DPO.
    \item Moreover, we investigate various synthetic answer selection strategies that rank candidates by the probability a panel of judges assigns to each answer's intended score level. We find that a mixture of answers where judges assign the requested label with confidence and uncertainty is the most useful for ASAS model training.
\end{itemize}
